\begin{frame}[plain]
\begin{tikzpicture}[remember picture,overlay]
\node at ([yshift=-80bp]current page.north) {\includegraphics[height=78bp]{\EMroot/ferdowsi-logo.png}};
\node[fill=SLbar,draw=SLbarEdge,line width=0.8bp,rounded corners=4bp,minimum width=640bp,minimum height=190bp] (b) at ([yshift=-262bp]current page.north) {};
\node[text=SLtitle,font=\bfseries\fontsize{40}{48}\selectfont] at ([yshift=52bp]b.center) {\EMdir{ریاضی مهندسی}};
\node[text=SLdef,font=\fontsize{20}{26}\selectfont,text width=600bp,align=center] at ([yshift=-4bp]b.center) {\EMdir{تحلیل فوریه \ \lr{\textbullet}\ معادلات دیفرانسیل با مشتقات جزئی \ \lr{\textbullet}\ آنالیز مختلط}};
\node[text=EMmuted,font=\fontsize{18}{22}\selectfont] at ([yshift=-58bp]b.center) {\EMdir{\lr{Engineering Mathematics}}};
\node[text=EMink,font=\fontsize{16}{20}\selectfont] at ([yshift=-410bp]current page.north) {\EMdir{استاد: عباس ابراهیمی مقدم \quad\lr{\textbullet}\quad جمع‌آوری و تنظیم: آیدین شکاری}};
\node[text=EMmuted,font=\fontsize{14}{18}\selectfont] at ([yshift=-440bp]current page.north) {\EMdir{دانشگاه فردوسی مشهد}};
\end{tikzpicture}
\end{frame}

\def\EMtitlesize{26}\def\EMtitleskip{31}
\begin{frame}{اطلاعات کلی درس}
\begin{uncoverenv}<2->
\EMp{\textbf{درس} \lr{—} ریاضی مهندسی (۲۳۱۱۲۰۱۸)}
\end{uncoverenv}
\EMgap
\begin{uncoverenv}<3->
\EMp{\textbf{مدرس} \lr{—} عباس ابراهیمی مقدم}
\end{uncoverenv}
\EMgap
\begin{uncoverenv}<4->
\EMp{\textbf{پیش‌نیاز درس} \lr{—} ریاضی ۲، معادلات دیفرانسیل}
\end{uncoverenv}
\EMgap
\begin{uncoverenv}<5->
\EMp{\textbf{امکانات آنلاین} \lr{—} پیام‌رسان پرتال، کانال تلگرام و نیز سایت \lr{vu}}
\end{uncoverenv}
\end{frame}

\def\EMtitlesize{26}\def\EMtitleskip{31}
\begin{frame}{منابع}
\begin{uncoverenv}<2->
\EMp{\textbf{منبع اصلی} \lr{—} نسخۀ الکترونیکی کتاب در سایت \lr{vu} بارگذاری شده است: \lr{\textit{Advanced Engineering Mathematics}, Erwin Kreyszig}}
\end{uncoverenv}
\EMgap
\begin{uncoverenv}<3->
\EMp{\textbf{سایر منابع}}
\end{uncoverenv}
\EMgap
\begin{uncoverenv}<4->
\EMitem{\lr{\textbullet}}{\dimexpr\linewidth-26bp\relax}{\lr{\textit{Advanced Engineering Mathematics}, Clarence R. Wylie}}
\end{uncoverenv}
\EMgap
\begin{uncoverenv}<5->
\EMitem{\lr{\textbullet}}{\dimexpr\linewidth-26bp\relax}{\lr{\textit{Complex Variables and Applications}, Ruel V. Churchill, \lr{\dots}}}
\end{uncoverenv}
\EMgap
\begin{uncoverenv}<6->
\EMitem{\lr{\textbullet}}{\dimexpr\linewidth-26bp\relax}{ریاضیات مهندسی، دکتر جلیل راشد محصل، دانشگاه تهران}
\end{uncoverenv}
\end{frame}

\def\EMtitlesize{26}\def\EMtitleskip{31}
\begin{frame}{نحوۀ ارزیابی درس}
\begin{uncoverenv}<2->
\EMitem{\lr{\textbullet}}{\dimexpr\linewidth-26bp\relax}{\textbf{۲۰\lr{٪}} \lr{—} کوییزها (در کلاس حل تمرین)}
\end{uncoverenv}
\EMgap
\begin{uncoverenv}<3->
\EMitem{\lr{\textbullet}}{\dimexpr\linewidth-26bp\relax}{\textbf{۱۰\lr{٪}} \lr{—} تکالیف}
\end{uncoverenv}
\EMgap
\begin{uncoverenv}<4->
\EMitem{\lr{\textbullet}}{\dimexpr\linewidth-26bp\relax}{\textbf{۳۰\lr{٪}} \lr{—} میان‌ترم (پنجشنبه ۲۵ آبان)}
\end{uncoverenv}
\EMgap
\begin{uncoverenv}<5->
\EMitem{\lr{\textbullet}}{\dimexpr\linewidth-26bp\relax}{\textbf{۴۰\lr{٪}} \lr{—} پایان‌ترم (روز سوم امتحانات)}
\end{uncoverenv}
\end{frame}

\def\EMtitlesize{26}\def\EMtitleskip{31}
\begin{frame}{موضوعات مطرح‌شده در درس}
\begin{uncoverenv}<2->
\EMp{\textbf{تحلیل فوریه}}
\end{uncoverenv}
\EMgap
\begin{uncoverenv}<3->
\EMitem{\lr{\textbullet}}{\dimexpr\linewidth-26bp\relax}{سری فوریۀ توابع متناوب}
\end{uncoverenv}
\EMgap
\begin{uncoverenv}<4->
\EMitem{\lr{\textbullet}}{\dimexpr\linewidth-26bp\relax}{انتگرال فوریۀ توابع غیرتناوبی}
\end{uncoverenv}
\EMgap
\begin{uncoverenv}<5->
\EMp{\textbf{معادلات دیفرانسیل با مشتقات جزئی}}
\end{uncoverenv}
\EMgap
\begin{uncoverenv}<6->
\EMitem{\lr{\textbullet}}{\dimexpr\linewidth-26bp\relax}{معرفی بعضی معادلات دیفرانسیل مهم}
\end{uncoverenv}
\EMgap
\begin{uncoverenv}<7->
\EMitem{\lr{\textbullet}}{\dimexpr\linewidth-26bp\relax}{روش‌های مختلف حل معادلات دیفرانسیل با مشتقات جزئی}
\end{uncoverenv}
\EMgap
\begin{uncoverenv}<8->
\EMp{\textbf{آنالیز مختلط}}
\end{uncoverenv}
\EMgap
\begin{uncoverenv}<9->
\EMitem{\lr{\textbullet}}{\dimexpr\linewidth-26bp\relax}{اعداد مختلط (روش‌های مختلف بیان اعداد مختلط، \lr{\dots})}
\end{uncoverenv}
\EMgap
\begin{uncoverenv}<10->
\EMitem{\lr{\textbullet}}{\dimexpr\linewidth-26bp\relax}{توابع مختلط (حد، پیوستگی، مشتق و انتگرال مختلط)}
\end{uncoverenv}
\EMgap
\begin{uncoverenv}<11->
\EMitem{\lr{\textbullet}}{\dimexpr\linewidth-26bp\relax}{سری لورنت}
\end{uncoverenv}
\end{frame}

\def\EMtitlesize{26}\def\EMtitleskip{31}
\begin{frame}{یادآوری بعضی نکات ریاضی}
\begin{uncoverenv}<2->
\EMitem{\lr{\textbullet}}{\dimexpr\linewidth-26bp\relax}{عملیات روی متغیر مستقل یک تابع: جمع متغیر مستقل با یک عدد ثابت، ضرب متغیر مستقل در یک عدد ثابت، قرینه کردن متغیر مستقل}
\end{uncoverenv}
\EMgap
\begin{uncoverenv}<3->
\EMitem{\lr{\textbullet}}{\dimexpr\linewidth-26bp\relax}{توابع زوج یا فرد و قسمت زوج یا فرد یک تابع}
\end{uncoverenv}
\EMgap
\begin{uncoverenv}<4->
\EMitem{\lr{\textbullet}}{\dimexpr\linewidth-26bp\relax}{توابع متناوبی و دورۀ تناوب این توابع}
\end{uncoverenv}
\end{frame}

\def\EMtitlesize{26}\def\EMtitleskip{31}
\begin{frame}{برگۀ فرمول: مشتق و انتگرال \lr{—} مشتق‌گیری}
\begin{uncoverenv}<2->
\EMeq{(cu)' = cu' \quad (c\ \EMtx{constant}{ثابت})}
\vspace{3bp}
\EMeq{(u+v)' = u'+v'}
\vspace{3bp}
\EMeq{(uv)' = u'v+uv'}
\vspace{3bp}
\EMeq{\left(\dfrac uv\right)' = \dfrac{u'v-uv'}{v^2}}
\vspace{3bp}
\EMeq{\dfrac{du}{dx} = \dfrac{du}{dy}\cdot\dfrac{dy}{dx}\quad(\EMtx{Chain rule}{قاعدۀ زنجیری})}
\end{uncoverenv}
\EMgap
\begin{uncoverenv}<3->
\EMeq{(x^n)' = nx^{n-1}}
\vspace{3bp}
\EMeq{(e^x)' = e^x}
\vspace{3bp}
\EMeq{(e^{ax})' = ae^{ax}}
\vspace{3bp}
\EMeq{(a^x)' = a^x\ln a}
\end{uncoverenv}
\end{frame}

\def\EMtitlesize{26}\def\EMtitleskip{31}
\begin{frame}{برگۀ فرمول: مشتق و انتگرال \lr{—} مشتق‌گیری (ادامه)}
\begin{uncoverenv}<2->
\EMeq{(\sin x)' = \cos x}
\vspace{3bp}
\EMeq{(\cos x)' = -\sin x}
\vspace{3bp}
\EMeq{(\tan x)' = \sec^2x}
\vspace{3bp}
\EMeq{(\cot x)' = -\csc^2x}
\vspace{3bp}
\EMeq{(\sinh x)' = \cosh x}
\vspace{3bp}
\EMeq{(\cosh x)' = \sinh x}
\vspace{3bp}
\EMeq{(\ln x)' = \dfrac1x}
\vspace{3bp}
\EMeq{(\log_a x)' = \dfrac{\log_a e}{x}}
\vspace{3bp}
\EMeq{(\arcsin x)' = \dfrac1{\sqrt{1-x^2}}}
\end{uncoverenv}
\end{frame}

\def\EMtitlesize{26}\def\EMtitleskip{31}
\begin{frame}{برگۀ فرمول: مشتق و انتگرال \lr{—} مشتق‌گیری (ادامه)}
\begin{uncoverenv}<2->
\EMeq{(\arccos x)' = -\dfrac1{\sqrt{1-x^2}}}
\vspace{3bp}
\EMeq{(\arctan x)' = \dfrac1{1+x^2}}
\vspace{3bp}
\EMeq{(\operatorname{arccot}x)' = -\dfrac1{1+x^2}}
\end{uncoverenv}
\end{frame}

\def\EMtitlesize{26}\def\EMtitleskip{31}
\begin{frame}{برگۀ فرمول: مشتق و انتگرال \lr{—} انتگرال‌گیری}
\begin{uncoverenv}<2->
\EMeq{\int uv'\,dx = uv-\int u'v\,dx \quad(\EMtx{by parts}{جزء‌به‌جزء})}
\vspace{3bp}
\EMeq{\int x^n\,dx = \dfrac{x^{n+1}}{n+1}+c \quad(n\neq-1)}
\vspace{3bp}
\EMeq{\int\dfrac1x\,dx = \ln\lvert x\rvert+c}
\vspace{3bp}
\EMeq{\int e^{ax}\,dx = \dfrac1a\,e^{ax}+c}
\end{uncoverenv}
\end{frame}

\def\EMtitlesize{26}\def\EMtitleskip{31}
\begin{frame}{برگۀ فرمول: مشتق و انتگرال \lr{—} انتگرال‌گیری (ادامه)}
\begin{uncoverenv}<2->
\EMeq{\int\sin x\,dx = -\cos x+c}
\vspace{3bp}
\EMeq{\int\cos x\,dx = \sin x+c}
\vspace{3bp}
\EMeq{\int\tan x\,dx = -\ln\lvert\cos x\rvert+c}
\vspace{3bp}
\EMeq{\int\cot x\,dx = \ln\lvert\sin x\rvert+c}
\vspace{3bp}
\EMeq{\int\sec x\,dx = \ln\lvert\sec x+\tan x\rvert+c}
\vspace{3bp}
\EMeq{\int\csc x\,dx = \ln\lvert\csc x-\cot x\rvert+c}
\end{uncoverenv}
\end{frame}

\def\EMtitlesize{26}\def\EMtitleskip{31}
\begin{frame}{برگۀ فرمول: مشتق و انتگرال \lr{—} انتگرال‌گیری (ادامه)}
\begin{uncoverenv}<2->
\EMeq{\int\dfrac{dx}{x^2+a^2} = \dfrac1a\arctan\dfrac xa+c}
\vspace{3bp}
\EMeq{\int\dfrac{dx}{\sqrt{a^2-x^2}} = \arcsin\dfrac xa+c}
\vspace{3bp}
\EMeq{\int\dfrac{dx}{\sqrt{x^2+a^2}} = \operatorname{arcsinh}\dfrac xa+c}
\vspace{3bp}
\EMeq{\int\dfrac{dx}{\sqrt{x^2-a^2}} = \operatorname{arccosh}\dfrac xa+c}
\end{uncoverenv}
\end{frame}

\def\EMtitlesize{26}\def\EMtitleskip{31}
\begin{frame}{برگۀ فرمول: مشتق و انتگرال \lr{—} انتگرال‌گیری (ادامه)}
\begin{uncoverenv}<2->
\EMeq{\int\sin^2x\,dx = \tfrac12x-\tfrac14\sin2x+c}
\vspace{3bp}
\EMeq{\int\cos^2x\,dx = \tfrac12x+\tfrac14\sin2x+c}
\vspace{3bp}
\EMeq{\int\tan^2x\,dx = \tan x-x+c}
\vspace{3bp}
\EMeq{\int\cot^2x\,dx = -\cot x-x+c}
\vspace{3bp}
\EMeq{\int\ln x\,dx = x\ln x-x+c}
\end{uncoverenv}
\EMgap
\begin{uncoverenv}<3->
\EMeq{\int e^{ax}\sin bx\,dx = \dfrac{e^{ax}}{a^2+b^2}\,(a\sin bx-b\cos bx)+c}
\vspace{3bp}
\EMeq{\int e^{ax}\cos bx\,dx = \dfrac{e^{ax}}{a^2+b^2}\,(a\cos bx+b\sin bx)+c}
\end{uncoverenv}
\end{frame}
